% 单文件讲义模板。当前正文仅为排版样页，新项目请替换为实际课程内容。
% 使用 XeLaTeX 编译；项目参数集中在本文件开头。
\documentclass[UTF8,fontset=none,zihao=-4,a4paper]{ctexart}

% ===== 可自主编辑的参数 =====
\newcommand{\LectureTitle}{学科讲义}
\newcommand{\LectureMainColor}{35566E} % 十六进制颜色，不写 #
\newcommand{\LectureBackColor}{F2F6F8}
\newcommand{\LectureBorderColor}{B8C8D3}
\newcommand{\LectureCJKFont}{Songti SC}
\newcommand{\LectureLatinFont}{Times New Roman}
\newcommand{\LectureCodeFont}{Menlo}
\newcommand{\LectureMathFont}{texgyretermes-math.otf}
\newcommand{\LectureBodySize}{12pt}
\newcommand{\LectureBodyBaseline}{14.4pt} % 默认字号的1.2倍
\newcommand{\LectureLineSpread}{1.22}
\newcommand{\LectureIndent}{2em}
\newcommand{\LectureParagraphSkip}{3pt}
\newcommand{\LectureLeftMargin}{26mm}
\newcommand{\LectureRightMargin}{26mm}
\newcommand{\LectureTopMargin}{25mm}
\newcommand{\LectureBottomMargin}{25mm}
\newcommand{\LectureSectionSize}{14pt}
\newcommand{\LectureSubsectionSize}{12pt}
\newcommand{\LectureCodeSize}{10pt}
\newcommand{\LectureTOCDepth}{1}
\newcommand{\LectureBookmarkDepth}{2}
\newcommand{\LectureBoxPadding}{3mm}
\newcommand{\LectureBoxRule}{0.45pt}
\newcommand{\LectureCoverSize}{36pt}
\newcommand{\LectureCoverOffset}{0.28\textheight}

% ===== 字体、主题和基础排版 =====
\usepackage[left=\LectureLeftMargin,right=\LectureRightMargin,
 top=\LectureTopMargin,bottom=\LectureBottomMargin,headheight=16pt]{geometry}
\IfFontExistsTF{\LectureCJKFont}{
 \setCJKmainfont{\LectureCJKFont}[AutoFakeBold=2]
 \setCJKsansfont{\LectureCJKFont}[AutoFakeBold=2]
 \setCJKmonofont{\LectureCJKFont}[AutoFakeBold=2]
}{
 \setCJKmainfont{FandolSong-Regular.otf}[BoldFont={FandolSong-Bold.otf}]
 \setCJKsansfont{FandolSong-Regular.otf}[BoldFont={FandolSong-Bold.otf}]
 \setCJKmonofont{FandolSong-Regular.otf}[BoldFont={FandolSong-Bold.otf}]
 \typeout{LECTURE-PREVIEW: CJK font replaced by FandolSong.}
}
\IfFontExistsTF{\LectureLatinFont}{\setmainfont{\LectureLatinFont}}{
 \setmainfont{texgyretermes-regular.otf}[BoldFont={texgyretermes-bold.otf},
 ItalicFont={texgyretermes-italic.otf},BoldItalicFont={texgyretermes-bolditalic.otf}]
 \typeout{LECTURE-PREVIEW: Latin font replaced by TeX Gyre Termes.}
}
\IfFontExistsTF{\LectureCodeFont}{\setmonofont{\LectureCodeFont}}{
 \setmonofont{lmmono10-regular.otf}
 \typeout{LECTURE-PREVIEW: Code font replaced by Latin Modern Mono.}
}
\usepackage{amsmath,amssymb,unicode-math}
\setmathfont{\LectureMathFont}
\usepackage{xcolor,booktabs,tabularx,array,enumitem,fancyhdr,needspace}
\usepackage[breakable,skins]{tcolorbox}
\usepackage{fvextra,tikz}
\usetikzlibrary{arrows.meta}
\usepackage[hidelinks]{hyperref}
\hypersetup{pdftitle={\LectureTitle},bookmarksdepth=\LectureBookmarkDepth}
\definecolor{ThemeMain}{HTML}{\LectureMainColor}
\definecolor{ThemeBack}{HTML}{\LectureBackColor}
\definecolor{ThemeBorder}{HTML}{\LectureBorderColor}
\newtcolorbox{infobox}[1]{enhanced,breakable,title={#1},
 colback=ThemeBack,colframe=ThemeBorder,colbacktitle=ThemeBack,
 coltitle=ThemeMain,fonttitle=\bfseries,
 attach title to upper={\par\smallskip},boxrule=\LectureBoxRule,
 arc=0.8mm,left=\LectureBoxPadding,right=\LectureBoxPadding,
 top=2mm,bottom=2mm,before skip=8pt,after skip=8pt,pad at break*=1.5mm}
\renewcommand{\normalsize}{\fontsize{\LectureBodySize}{\LectureBodyBaseline}\selectfont}
\linespread{\LectureLineSpread}
\ctexset{section={format=\fontsize{\LectureSectionSize}{18pt}\selectfont\bfseries,
 beforeskip=0pt,afterskip=14pt},
 subsection={format=\fontsize{\LectureSubsectionSize}{15pt}\selectfont\bfseries,
 beforeskip=12pt,afterskip=6pt}}
\AtBeginDocument{\normalsize\setlength{\parindent}{\LectureIndent}}
\setlength{\parskip}{\LectureParagraphSkip}
\setcounter{tocdepth}{\LectureTOCDepth}
\setcounter{secnumdepth}{2}
\setlength{\tabcolsep}{5pt}
\renewcommand{\arraystretch}{1.3}
\setlist{itemsep=3pt,topsep=4pt,parsep=0pt,leftmargin=2em}
\newcolumntype{Y}{>{\raggedright\arraybackslash}X}
\setlength{\emergencystretch}{1.5em}
\allowdisplaybreaks[2]\raggedbottom
\pagestyle{fancy}\fancyhf{}
\fancyhead[L]{\small\LectureTitle}
\fancyhead[R]{\small\nouppercase{\leftmark}}
\fancyfoot[C]{\small\thepage}
\renewcommand{\headrulewidth}{0.3pt}
\renewcommand{\sectionmark}[1]{\markboth{#1}{}}
\newcommand{\chapterpage}[1]{\clearpage\section{#1}}
\newcommand{\term}[1]{\textbf{#1}}
\newcommand{\termdef}[2]{\hypertarget{term:#1}{\textbf{#2}}}
\newcommand{\termref}[2]{\hyperlink{term:#1}{#2}}
\newcommand{\exercise}[2]{\hypertarget{exercise:#1}{\textbf{#2}}\quad\hyperlink{answer:#1}{\small 查看答案}\quad}
\newcommand{\answer}[2]{\hypertarget{answer:#1}{\textbf{#2}}\quad\hyperlink{exercise:#1}{\small 返回题目}\quad}
\newcommand{\sol}{\par\noindent\textbf{解}\quad}
\numberwithin{equation}{section}
\DefineVerbatimEnvironment{codeblock}{Verbatim}{
 fontsize=\fontsize{\LectureCodeSize}{12pt}\selectfont,
 breaklines=true,breakanywhere=false,frame=single,
 rulecolor=\color{ThemeBorder},framesep=2mm,numbers=left,numbersep=6pt,samepage=true}

\begin{document}
\begin{titlepage}\thispagestyle{empty}\vspace*{\LectureCoverOffset}
\begin{center}{\fontsize{\LectureCoverSize}{45pt}\selectfont\bfseries\LectureTitle}\end{center}
\end{titlepage}
\pagenumbering{Roman}\thispagestyle{empty}\tableofcontents
\clearpage\pagenumbering{arabic}

\section{排版样页}\label{ch:sample}
\begin{infobox}{样页用途}
本页展示正文、公式、表格、图示和信息框。以下数值为固定教学示例，仅供检查排版；正式写作时，应换成该课程的资料与记号。
\end{infobox}

\subsection{正文与公式}
\termdef{mean}{算术平均数}是各个数值之和除以数值个数。设有$n$个观测值$x_1,\ldots,x_n$，其中$n$为正整数，记平均数为$\bar{x}$，则
\begin{equation}\bar{x}=\frac{1}{n}\sum_{i=1}^{n}x_i.\end{equation}
这里的下标$i$表示观测次序。若观测值以厘米计，平均数也以厘米计。实际课程中应沿用用户资料的符号，不因模板使用了$x$就将原资料中的其他字母替换为$x$。

\begin{infobox}{例题1.1：平均数计算}
三个观测值为$2,4,6$，求其算术平均数。
\end{infobox}
\sol 数值个数为$n=3$，总和为$2+4+6=12$，因此$\bar{x}=12/3=4$。

\subsection{表格与图示}
\begin{center}
\begin{tabularx}{\linewidth}{c c Y}
\toprule
观测次序 & 观测值 & 说明\\\midrule
1 & 2 & 保留输入顺序\\
2 & 4 & 此值恰好等于平均数\\
3 & 6 & 样例不代表实际测量结果\\
\bottomrule
\end{tabularx}
\end{center}

\begin{center}
\begin{tikzpicture}[font=\small,>=Stealth]
\draw[->,color=ThemeMain] (0,0)--(6.5,0) node[right]{$i$};
\draw[->,color=ThemeMain] (0,0)--(0,2.8) node[above]{$x_i$};
\draw[dashed,color=ThemeBorder] (0,1.6)--(6,1.6);
\node[anchor=west] at (6.05,1.6) {$\bar{x}=4$};
\foreach \pos/\val/\height in {1/2/0.8,3/4/1.6,5/6/2.4}{
 \fill[color=ThemeMain] (\pos,\height) circle (2pt);
 \node[below] at (\pos,0) {\pgfmathparse{(\pos+1)/2}\pgfmathprintnumber{\pgfmathresult}};
 \node[above] at (\pos,\height) {\val};
}
\end{tikzpicture}
\end{center}
横轴为观测次序，纵轴为观测值；虚线表示平均数。图中文字、线条和主题色应与正文协调。

\newpage
\subsection{代码、练习与答案}
这段Python代码与前页的计算对应，输入固定为三个数值。非编程课程可以移除此样例，改为材料分析或实验步骤。
\begin{codeblock}
values = [2, 4, 6]
mean = sum(values) / len(values)
print(mean)  # 4.0
\end{codeblock}
\begin{infobox}{注意：教学条件}
上述计算要求输入非空。这里的简单样例不代替面向任意输入的完整程序；专业讲义应说明自身的简化条件。
\end{infobox}
\subsection{练习}
\par\noindent\exercise{mean-1}{1.1}将观测值改为$2,4,9$，计算新的平均数并解释其变化。
\par\noindent\exercise{mean-2}{1.2}说明平均数为什么具有与原观测值相同的单位。
\subsection{参考答案}
\par\noindent\answer{mean-1}{1.1}总和为15，观测数仍为3，因此平均数为5。第三个数值增加3，总和也增加3；除以3后，平均数增加1。

\par\noindent\answer{mean-2}{1.2}相加的观测值单位相同，除以无单位的计数$n$后，结果仍具有原单位。
\subsection{术语索引}
\noindent\termref{mean}{算术平均数}\quad\hyperref[ch:sample]{本章定义与计算示例}
\end{document}
